\chapter{Building and Measuring an Optimizing Compiler}\label{ch:building-measuring-optimizing-compiler}

A collection of passes is not yet a compiler.
Passes must run in an order that respects their contracts, share facts without corruption, and stop within budget.
Then the result must be measured honestly, because an unverified speedup claim is a defect in the engineering record.
This chapter assembles the pipeline and the measurement discipline together.
We specify pass ordering, fact management, and a reproducible checklist that reports conditions instead of inventing numbers.

\section{From Passes to a Pipeline}\label{sec:ch38-pipeline}

A pipeline is a sequence of passes joined by contracts.
Each pass declares what it needs, what it preserves, and what it establishes.
The driver orders passes so every need is established earlier and no later pass silently breaks an earlier guarantee.
Figure~\ref{fig:ch38-pipeline} shows the canonical shape used in this book.
Front-end passes build a valid intermediate form.
Middle-end passes analyze and transform in alternating layers.
Back-end passes lower, allocate, and emit with machine facts checked at each step.
Table~\ref{tab:ch38-ordering} lists common ordering rules with the defect each rule prevents.
Analysis before transformation prevents stale facts from guiding rewrites.
Cleanup before measurement prevents dead code from timing the benchmark.
Lowering late prevents machine details from blocking portable optimization.
Each rule exists because violating it once produced a real wrong-code or slow-code outcome.
\term{Phase ordering} is the problem of choosing the sequence; \term{fact invalidation} is the discipline of dropping facts a transformation may have broken.
A pass that moves stores must drop alias facts it cannot re-derive.
A pass that splits edges must update dominance information or mark it stale.
The safest policy recomputes cheap facts and versions expensive ones.
Every fact carries the pass and revision that produced it.
Any consumer checks the version before trusting the content.

\begin{figure}[htbp]
\centering
\begin{tikzpicture}[node distance=0.9cm and 1.2cm]
\node[basic block] (fe) {front end\\ valid IR};
\node[basic block, right=of fe] (mid) {middle end\\ analyze, transform};
\node[basic block, right=of mid] (be) {back end\\ lower, allocate, emit};
\node[basic block, below=of mid] (check) {contract checks\\ at each handoff};
\draw[cfg edge] (fe) -- (mid);
\draw[cfg edge] (mid) -- (be);
\draw[cfg edge] (mid) -- (check);
\draw[cfg edge] (check) -- (mid);
\end{tikzpicture}
\caption{Passes flow from language form through optimization to machine form with contract checks at each handoff, so ordering defects surface where they occur.}
\label{fig:ch38-pipeline}
\end{figure}

\begin{table}[htbp]
\centering
\caption{Ordering rules pair each sequence constraint with the defect it prevents.}
\label{tab:ch38-ordering}
\begin{tabular}{lll}
\toprule
Rule & Order & Prevents \\
\midrule
Analyze then change & Facts before rewrites & Stale-fact miscompilation \\
Clean then measure & Cleanup before timing & Dead-code timing noise \\
Lower late & Machine detail last & Blocked portable passes \\
\bottomrule
\end{tabular}
\end{table}

\section{Managing Facts Across Passes}\label{sec:ch38-facts}

Facts flow between passes through a shared structure with explicit ownership.
Each analysis publishes its result with a version stamp.
Each transformation declares which stamps it preserves, updates, or destroys.
Algorithm~\ref{alg:ch38-pipeline-run} gives the driver loop that enforces this discipline.
It assumes every pass declares its needs, preserves, and destroys sets over named facts.
It requires contract checks at each handoff, reusing the gates of the previous chapter.
It runs passes in order, recomputes missing facts on demand, and stops at the first contract violation with the suspect named.
Preservation claims are checked, not trusted.
A pass that claims to preserve dominance but moves a store through a join faces the validator before the next pass runs.
Budgets bound the whole run.
An optimization budget caps total compile time and code growth.
A pass that exceeds its share yields to the remaining pipeline rather than starving it.
Determinism is a further requirement.
The same input with the same flags and the same environment must produce the same output.
Nondeterministic iteration orders, unseeded hashes, and parallel pass schedules all threaten that property.
We fix iteration orders, seed every random choice, and log the seeds with the build record.

\begin{algorithm}[htbp]
\caption{Contract-checked pipeline driver with fact versions}\label{alg:ch38-pipeline-run}
\begin{pseudocode}
\Require pass list $P_1 \dots P_k$, input program $T$, fact store $S$
\Ensure output passing the enabled checks, or the first detected contract violation
\ForAll{$i$ from $1$ to $k$}
\State check entry contract of $P_i$ on $T$ and $S$
\If{entry check fails}
\State \Return violation at entry of $P_i$
\EndIf
\State run $P_i$ on $T$ and update $S$ \Comment{drop destroyed facts}
\State check exit contract of $P_i$ on $T$ and $S$
\If{exit check fails}
\State \Return violation at exit of $P_i$
\EndIf
\EndFor
\State \Return $T$ with build record
\end{pseudocode}
\end{algorithm}

\section{Worked Example: Ordering Three Passes}\label{sec:ch38-worked}

We order folding, dead-code elimination, and register allocation for a small body.
The body computes \code{t1 = 3 + 5}, \code{t2 = t1 + x}, and \code{t3 = t1 + 0}, then returns \code{t2}.
Folding runs first because it needs only well-formed trees and establishes new constants.
Hand evaluation gives $t1 = 8$ in every execution.
Folding without operand substitution leaves \code{t1 = 8}, \code{t2 = t1 + x}, and \code{t3 = t1 + 0}, then returns \code{t2}.
Dead-code elimination removes \code{t3} because its result is unused; that fact does not require constant folding.
It keeps \code{t1} because \code{t2} still uses it.
If constant propagation instead substitutes $8$ into \code{t2}, the now-unused \code{t1} can also be deleted.
Register allocation runs last because it needs the final live ranges.
Live analysis gives \code{t1} live across one use and \code{t2} live to the return.
Two registers or one register plus forwarding suffice; the exact choice is target-specific.
Reversing the first two passes removes \code{t3} just as safely in this example.
Folding can expose additional dead paths in other programs, which motivates repeating cleanup after it.
Running allocation first spends effort on an unused definition that cleanup could have removed earlier.
The listing below pins the final body for tests.
\begin{irlisting}
; after fold, eliminate, allocate: order matters
t1 = 8            ; folded from 3 + 5
t2 = add t1, x    ; kept, used by return
ret t2            ; t3 removed as dead
\end{irlisting}
Each step preserves the return value $8+x$ under the IR's arithmetic rules.
A differential test should compare that result against the original program; this is a hand-worked example, not a report of an implemented pipeline.

\section{A Reproducible Measurement Checklist}\label{sec:ch38-checklist}

Performance claims require a measurement method that others can repeat.
This section gives a checklist, not numbers.
No timing figures appear here because none were collected for this chapter.
Any claim about speedup needs a fresh run recorded under the headings below.
The checklist has three phases: preparation, execution, and checking.
Preparation fixes everything except the variable under test.
Execution collects observations under stated conditions.
Checking validates the observations before any conclusion is drawn.
\begin{ccnote}[Report conditions]
A measurement without its conditions is an anecdote.
Record the full environment with every result so others can reproduce or refute it.
\end{ccnote}
Preparation steps come first.
Select the benchmark programs and inputs, and record their versions and sources.
Fix the compiler flags, the optimization levels, and the complete build record.
Fix the machine, operating system, and library versions under test.
Disable or document background load, frequency scaling, and competing processes.
Pre-build all binaries and verify checksums before timing begins.
Decide the warm-up policy for just-in-time systems and document steady-state criteria.
Execution steps come next.
Run each configuration repeatedly with a fixed repetition count chosen in advance.
Randomize run order across configurations to spread interference evenly.
Record wall time, CPU time, and any allocator or collector statistics the platform exposes.
Record failures, timeouts, and discarded runs with their reasons rather than silently dropping them.
Keep raw logs; never keep only averages.
Checking steps come last.
Verify that contract checks passed on every measured build.
Verify checksums of binaries and inputs before accepting any observation.
Summarize with distributions and counts, not single best runs.
Report the machine, flags, inputs, repetitions, warm-up, and raw-data location alongside every figure.
State threats explicitly: what could still differ between configurations and what was done about each item in Table~\ref{tab:ch38-threats}.
The \enquote{conditions not numbers} rule keeps the record honest.
It is always acceptable to report that an experiment has not been run.
It is never acceptable to present an expected speedup as a measured one.

\begin{table}[htbp]
\centering
\caption{Common measurement threats pair each risk with its mitigation in the checklist.}
\label{tab:ch38-threats}
\begin{tabular}{ll}
\toprule
Threat & Mitigation \\
\midrule
Background load & Document and isolate machine \\
Frequency scaling & Fix or record governor policy \\
Cold caches & Documented warm-up excluded \\
Single best run & Fixed repetitions, full logs \\
\bottomrule
\end{tabular}
\end{table}

\section{Implementation Notes and Complexity}\label{sec:ch38-implementation}

Driver cost is linear in the pass count plus the cost of contract checks.
If each enabled structural check is linear in the current IR size $N_i$, total check cost is $O(\sum_{i=1}^{k}N_i)$, or $O(kN)$ when every size is bounded by $N$.
Dominance recomputation, semantic equivalence checking, and other nonlocal validators need their own bounds; the driver does not make arbitrary contracts linear-time checkable.
Fact recomputation adds the cost of analyses rerun after destructive passes.
Version stamps keep that cost explicit: every dropped fact names the pass that dropped it.
Production drivers add pass groups, function-granular scheduling, and parallel execution with deterministic merge.
They cache analysis results across functions and invalidate by revision.
They record the full build description: pass list, order, flags, seeds, versions, and input checksums.
A common pitfall adds a pass without declaring its destroyed facts.
Later passes then trust stale information and miscompile quietly.
Another pitfall tunes the pipeline against one benchmark and reports only that result.
The checklist prevents the second failure by requiring named benchmarks, fixed repetitions, and stated threats.
We treat the build record as part of the output.
A binary without its record is unrepeatable and therefore unmeasurable.
Figure~\ref{fig:ch38-pipeline} and Algorithm~\ref{alg:ch38-pipeline-run} together define the minimum a driver must implement.
Everything else is tuning with evidence.
% TODO(cite): need a primary reference for pipeline-driver determinism practices.

\section{Exercises}\label{sec:ch38-exercises}

\begin{enumerate}
\item Order four passes (fold, inline, dead-code, allocate) for a call-heavy body, justify each position with Table~\ref{tab:ch38-ordering}, and name the facts each pass destroys.\label{ex:ch38-order-passes}
\item Write a build record for a hypothetical experiment you have not run, listing every checklist field with explicitly empty result slots and a statement that no numbers are claimed.\label{ex:ch38-record-empty}
\item Design a threat analysis for comparing two optimization levels on one machine, covering each row of Table~\ref{tab:ch38-threats} with a concrete mitigation you can perform.\label{ex:ch38-threat-plan}
\end{enumerate}

\section{Further Reading}\label{sec:ch38-reading}

Appel presents complete compiler pipelines from source to machine code~\cite{appel-1998-modern}.
Morgan discusses pass organization and its effect on optimization quality~\cite{morgan-1998-optimizing}.
Safonov treats measurement and validation as trustworthiness practices~\cite{safonov-2010-trustworthy}.
The C standard draft defines conformance and observable behavior for C-facing claims~\cite{wg14-2011-c11}.
